\section{Solid-reference form for porous media modeling}
\label{sec:reservoir_simulation}

This section rewrites the transport and mechanics equations in a form suitable
for solution on the trajectory of the solid skeleton.  The solid motion
$\bs\chi$ supplies the map from the solid reference coordinate
$\mbf X$ to the current coordinate $\mbf x$, with deformation gradient
$\mbf F$ and Jacobian $J$.

\subsection{Relative flux vectors}
\label{sec:relative_flux_vectors}

The phase momentum equation with dynamic pressure lag
\eqref{eq:el_mom_f_dynamic_capillary} determines
Darcy-scale relative fluxes once interphase drag is written in terms of a
velocity relative to the solid skeleton.  For each fluid phase \(f\), define the
current relative mass flux by
%
\begin{equation}
	\mbf w_f
	=\phi_f\bar\rho_f\left(\mbf v_f-\mbf v_{\mc S}\right)
	=\rho_f\left(\mbf v_f-\mbf v_{\mc S}\right),
	\qquad
	f\in\mc F .
	\label{eq:relative_mass_flux_definition}
\end{equation}
%
Thus
%
\begin{equation}
	\mbf v_f
	=\mbf v_{\mc S}+\frac{\mbf w_f}{\rho_f},
	\qquad
	f\in\mc F .
	\label{eq:relative_flux_velocity_reconstruction}
\end{equation}
%
A sufficient fluid--skeleton interaction closure consistent with
\eqref{eq:MC_interphase_exchange_law} is
%
\begin{equation}
	\mbf b_f
	=
	-\frac{\mu_f^{\mathrm{visc}}\phi_f}{\bar\rho_f}\,
	\bs\kappa_f^{-1}\mbf w_f,
	\qquad
	f\in\mc F .
	\label{eq:fluid_phase_interaction_law}
\end{equation}
%
Here $\mu_f^{\mathrm{visc}}$ is the dynamic viscosity and $\bs\kappa_f$ is the
permeability tensor for phase $f$.  Phase relative permeability may be absorbed
into $\bs\kappa_f$ and may use the same capillary-history state \(\mbf h\) when
relative-permeability hysteresis is retained.
Equation~\eqref{eq:fluid_phase_interaction_law} contains
only dissipative drag; all reversible driving forces remain in the phase
momentum balance.  From \eqref{eq:relative_mass_flux_definition}, the two drag
tensors are related by
\(\mbf k_{f\mc S}^{-1}
=\phi_f^2\mu_f^{\mathrm{visc}}\bs\kappa_f^{-1}\).
Thus positive definiteness of \(\bs\kappa_f\), together with the
interaction-energy allocation
\eqref{eq:MC_drag_heating_equal_partition}, makes
\eqref{eq:MC_mechanical_interaction_admissibility} nonnegative.

Substituting \eqref{eq:relative_flux_velocity_reconstruction} and
\eqref{eq:fluid_phase_interaction_law} into the fluid momentum balance
\eqref{eq:el_mom_f_dynamic_capillary} gives, for each fluid phase,
%
\begin{align}
	\frac{1}{\rho_f}
	\left(
		\phi_f^2\mu_f^{\mathrm{visc}}\bs\kappa_f^{-1}
		+\sum_\alpha\dot{c}_f^\alpha\mbf I
	\right)\mbf w_f
		= & \;
		\rho_f\left(\mbf g-\mbf a_f\right)
		-\phi_f\nabla_{\mbf x}
	\left(
		\bar p_f-\omega_f^+
	\right)
		\notag \\
	  &
		+\nabla_{\mbf x}\cdot
	\left(
		\phi_f\mbf E\otimes\mbf d_f
	\right)
		+\sum_\alpha\dot{c}_f^\alpha
	\left(\nabla_{\mbf x}\tau-\mbf v_{\mc S}\right),
	\qquad
	f\in\mc F,
	\label{eq:relative_flux_linear_system}
\end{align}
%
By the component mass balance and the separate zero-sum closure for the
component-relative fluxes, the scalar
\(\sum_\alpha\dot{c}_f^\alpha\) is the net phase-mass source due to
conversion.  It enters the phase momentum balance as a source coefficient
weighting the conversion-driven relative-flow term there.  The
component-relative transport itself is closed through
\eqref{eq:MC_relative_transport_closures}.

To solve the preceding linear system for \(\mbf w_f\), invert the coefficient
multiplying that flux.  Define the resulting conversion-corrected mobility
tensor by
%
\begin{equation}
	\widetilde{\bs\lambda}_f
	=
	\left(
		\phi_f^2\mu_f^{\mathrm{visc}}\bs\kappa_f^{-1}
		+\sum_\alpha\dot{c}_f^\alpha\mbf I
	\right)^{-1},
	\qquad
	f\in\mc F .
	\label{eq:modified_relative_flux_permeability}
\end{equation}
%
The resistance inside the inverse has units
\({\rm kg\,m^{-3}\,s^{-1}}\), and
\(\widetilde{\bs\lambda}_f\) therefore has units
\({\rm m^3\,s\,kg^{-1}}={\rm m^2\,Pa^{-1}\,s^{-1}}\).  The underlying
drag law is dissipative when the phase permeability
\(\bs\kappa_f\) is symmetric positive definite, \(\mu_f^{\mathrm{visc}}>0\), and
\(\phi_f>0\).  The added scalar conversion-insertion term comes from the
momentum source as the conversion-insertion contribution.  Because \(\bs\kappa_f\) is symmetric with eigenvalues
\(\kappa_{f,i}\), the combined resistance
\(\phi_f^2\mu_f^{\mathrm{visc}}\bs\kappa_f^{-1}
+\sum_\alpha\dot{c}_f^\alpha\mbf I\)
has eigenvalues
\(\lambda_{f,i}^{\mathrm{res}}=\phi_f^2\mu_f^{\mathrm{visc}}/\kappa_{f,i}
+\sum_\alpha\dot{c}_f^\alpha\),
for each \(f\in\mc F\).  These are eigenvalues of the combined resistance,
not of the mobility \(\widetilde{\bs\lambda}_f\), whose eigenvalues are the
reciprocals \(1/\lambda_{f,i}^{\mathrm{res}}\).  The entropy-inequality
admissibility conditions require these resistance eigenvalues to be positive,
so the combined resistance is positive definite and the mobility
\(\widetilde{\bs\lambda}_f\) is well
defined.  The binding eigenvalue is that in the direction of largest phase
permeability.  If \(\kappa_{f,\max}\) and \(\kappa_{f,\min}\) are the
largest and smallest eigenvalues of \(\bs\kappa_f\), respectively, the
admissibility condition is equivalently
%
\begin{equation}
	\sum_\alpha\dot{c}_f^\alpha
	>
	-\frac{\phi_f^2\mu_f^{\mathrm{visc}}}{\kappa_{f,\max}},
	\qquad
	f\in\mc F .
	\label{eq:modified_relative_flux_scalar_admissibility}
\end{equation}
%
Within this positive-definite regime, the spectral condition number is
%
\begin{equation}
	\operatorname{cond}_2
	\left(
		\phi_f^2\mu_f^{\mathrm{visc}}\bs\kappa_f^{-1}
		+\sum_\alpha\dot{c}_f^\alpha\mbf I
	\right)
	=
	\frac{
		\phi_f^2\mu_f^{\mathrm{visc}}/\kappa_{f,\min}
		+\sum_\alpha\dot{c}_f^\alpha
	}{
		\phi_f^2\mu_f^{\mathrm{visc}}/\kappa_{f,\max}
		+\sum_\alpha\dot{c}_f^\alpha
	},
	\qquad
	f\in\mc F .
	\label{eq:modified_relative_flux_condition_number}
\end{equation}
%
Thus quantitative conditioning requires the denominator to remain bounded away
from zero, not merely positive.
For an aqueous reaction that removes fluid mass on a time scale \(t_{\rm rxn}\),
\(\sum_\alpha\dot{c}_f^\alpha\simeq-\rho_f/t_{\rm rxn}\).  The singular threshold
at the first loss of positive definiteness is
%
\begin{equation}
	t_{\rm rxn}
	=
	\frac{\rho_f\kappa_{f,\max}}{\phi_f^2\mu_f^{\mathrm{visc}}},
	\qquad
	f\in\mc F .
	\label{eq:modified_relative_flux_reaction_timescale_threshold}
\end{equation}
%
For longer time scales the combined resistance remains positive definite.  At
equality it is singular in the largest-permeability direction; at shorter time
scales the admissibility condition fails.
With representative reservoir values
\(\rho_f=10^3\,{\rm kg\,m^{-3}}\), \(\mu_f^{\mathrm{visc}}=10^{-3}\,{\rm Pa\,s}\),
\(\phi_f=0.3\), and \(\kappa_{f,\max}=10^{-12}\,{\rm m^2}\), this time is about
\(10^{-5}\,{\rm s}\), corresponding to a removal rate of order
\(10^8\,{\rm kg\,m^{-3}\,s^{-1}}\).  For conversion rates whose characteristic
time is much longer than this value, the source correction is a small
perturbation of the Darcy drag.  The estimate supplies a lower reaction-time
bound for the positive-resistance form of the algebraic relative-flux closure.

The resulting relative mass flux is
%
\begin{empheq}[box=\fbox]{align}
	\mbf w_f
	= & \;
	\rho_f\widetilde{\bs\lambda}_f
	\left[
		\rho_f\left(\mbf g-\mbf a_f\right)
		-\phi_f\nabla_{\mbf x}
		\left(
			\bar p_f-\omega_f^+
		\right)
		\right.
		\notag \\
	& \phantom{\;\rho_f\widetilde{\bs\lambda}_f}
		\left.
		+\nabla_{\mbf x}\cdot
		\left(
			\phi_f\mbf E\otimes\mbf d_f
		\right)
		+\sum_\alpha\dot{c}_f^\alpha
		\left(
			\nabla_{\mbf x}\tau-\mbf v_{\mc S}
		\right)
		\right],
	\qquad
	f\in\mc F .
	\label{eq:generalized_fluid_darcy_mass_flux}
\end{empheq}
%
This expression is the multicomponent fluid-phase analogue of a generalized
Darcy flux.  Its driving force contains body force and inertia, the actual
phase-pressure gradient assembled from stored capillarity, electrical
corrections, and any dynamic lag, together with Maxwell stress and conversion
insertion.
The observer-conversion term, the relative-velocity contribution that arises
when the phase-following thermodynamic rate is expressed in the skeleton
frame, belongs to that single phase-following rate rather than entering the
flux as an independent driving force.
Equation~\eqref{eq:MC_collected_observer_correction_reassembly} recombines it
with the skeleton-rate free-energy contribution before the collected
inequality is separated into reversible and residual parts.
Chemical-potential gradients drive component-relative transport
through \eqref{eq:MC_relative_transport_closures}, a distinct dissipative
mechanism alongside the bulk phase force and the conversion variation of
\(\tau\).  The admissibility conditions above guarantee a positive-definite
combined resistance, so algebraic flux elimination yields a positive-resistance
operator; thermodynamic dissipation resides in the drag and
component-transport closures, and the interaction-energy allocation of
\eqref{eq:MC_drag_heating_equal_partition} is retained.

The neutral, conversion-free, negligible-inertia specialization
\(\omega_f^+=0\), \(\mbf E\otimes\mbf d_f=\mbf0\),
\(\dot{c}_f^\alpha=0\), and \(\mbf a_f=\mbf0\) reduces
\eqref{eq:generalized_fluid_darcy_mass_flux} to
%
\begin{equation}
	\phi_f\left(\mbf v_f-\mbf v_{\mc S}\right)
	=
	-\frac{\bs\kappa_f}{\mu_f^{\mathrm{visc}}}
	\left(
		\nabla_{\mbf x}\bar p_f
		-\bar\rho_f\mbf g
	\right),
	\qquad
	f\in\mc F .
	\label{eq:standard_darcy_mass_flux_limit}
\end{equation}
%
This is Darcy's law for phase \(f\) relative to the deforming solid skeleton.
Neither isothermality nor a spatially uniform Helmholtz energy is required.
Nonisothermal cross effects may modify this law only when they are
introduced explicitly through a coupled constitutive closure.

\subsection{Solid-reference relative fluxes}
\label{sec:solid_reference_relative_fluxes}

The current relative flux in \eqref{eq:generalized_fluid_darcy_mass_flux} is
converted to a flux through surfaces in the solid reference configuration by the
Piola transform.  Define the reference relative mass flux
%
\begin{equation}
	\mbf W_f
	=
	J\mbf F^{-1}\mbf w_f,
	\qquad
	f\in\mc F .
	\label{eq:reference_relative_mass_flux}
\end{equation}
%
Define the pulled-back conversion-corrected mobility tensor
%
\begin{equation}
	\widetilde{\bs\Lambda}_f
	=
	J\mbf F^{-1}
	\widetilde{\bs\lambda}_f
	\mbf F^{-T},
	\qquad
	f\in\mc F .
	\label{eq:pulled_back_modified_permeability}
\end{equation}
%
The flux assembled in the reference component balances is then
%
\begin{align}
	\mbf W_f
	= & \; \rho_f\widetilde{\bs\Lambda}_f
	\left[
		\rho_f\mbf F^T\left(\mbf g-\mbf a_f\right)
		-\phi_f\nabla_{\mbf X}
		\left(
			\bar p_f-\omega_f^+
		\right)
		\right.
		\notag \\
	& \phantom{\;\rho_f\widetilde{\bs\Lambda}_f}
		\left.
		+\mbf F^T
		\left[
			\nabla_{\mbf x}\cdot
			\left(
				\phi_f\mbf E\otimes\mbf d_f
			\right)
		\right]
		+\sum_\alpha\dot{c}_f^\alpha
	\left(\nabla_{\mbf X}\tau-\mbf F^T\mbf v_{\mc S}\right)
	\right],
	\qquad
	f\in\mc F .
	\label{eq:solid_reference_relative_flux}
\end{align}
%
Here gradients of scalar fields obey
$\nabla_{\mbf x}(\cdot)=\mbf F^{-T}\nabla_{\mbf X}(\cdot)$, while the
divergence of a spatial flux obeys
%
\begin{equation}
	J\nabla_{\mbf x}\cdot\mbf w_f
	=
	\nabla_{\mbf X}\cdot
	\mbf W_f,
	\qquad
	f\in\mc F .
	\label{eq:solid_reference_flux_divergence}
\end{equation}
%
Equations~\eqref{eq:reference_relative_mass_flux}--\eqref{eq:solid_reference_relative_flux}
and \eqref{eq:solid_reference_flux_divergence} give the reference flux used in
the multicomponent balances.  The expression retains phase labels, component
conversion rates, and gradients of \(\tau\).

\subsection{Component balances on the solid trajectory}
\label{sec:solid_reference_component_balances}

For a fluid phase, substituting
\eqref{eq:relative_flux_velocity_reconstruction} into
\eqref{eq:averaged_component_spatial_mass}, using
\eqref{eq:component_partial_density_representation}, and writing the storage
term on the solid trajectory gives
%
\begin{equation}
	\left.
	\frac{\partial}{\partial t}
	\left(J\phi_f\bar\rho_f\eta_f^\alpha\right)
	\right\vert_{\mbf X}
	+\nabla_{\mbf X}\cdot
	\left[
	\eta_f^\alpha\mbf W_f
	+
	J\mbf F^{-1}
	\left(
	\mbf j_{f,{\rm disp}}^\alpha
	+
	\mbf j_{f,{\rm diff}}^\alpha
	\right)
	\right]
	=
	J\sum_m\nu_{f (m)}^\alpha r_{(m)},
	\qquad
	f\in\mc F .
	\label{eq:solid_reference_fluid_component_balance}
\end{equation}
%
For a solid phase \(s\), the phase velocity is the skeleton velocity, so there
is no bulk phase flux relative to the solid mesh.  Component diffusion and
dispersion remain admissible, giving
%
\begin{equation}
	\begin{aligned}
		\left.
		\frac{\partial}{\partial t}
		\left(J\phi_s\bar\rho_s\eta_s^\alpha\right)
		\right\vert_{\mbf X}
		& {}+
		\nabla_{\mbf X}\cdot
		\left[
			J\mbf F^{-1}
			\left(
				\mbf j_{s,{\rm disp}}^\alpha
				+
				\mbf j_{s,{\rm diff}}^\alpha
			\right)
		\right]
		\\
		& {}=
			J\sum_m\nu_{s (m)}^\alpha r_{(m)},
		\qquad
		s\in\mc S .
	\end{aligned}
	\label{eq:solid_reference_solid_component_balance}
\end{equation}
%
Thus the solid-reference storage variables are the Lagrangian component masses
\(J\phi_\xi\bar\rho_\xi\eta_\xi^\alpha\).  Transport through the solid mesh
contains bulk fluid motion relative to the skeleton and component diffusion or
dispersion relative to the fluid or solid phase.

\subsection{Overall momentum on the solid trajectory}
\label{sec:solid_reference_overall_momentum}

Define the general solid-reference effective first Piola--Kirchhoff stress by
%
\begin{equation}
	\mbf P^{\prime\prime}
	\equiv
	J\bs\sigma^{\prime\prime}\mbf F^{-T}.
	\label{eq:solid_reference_effective_piola_stress}
\end{equation}
%
The overall momentum balance \eqref{eq:MC_overall_momentum_nonlinear_biot} can be
written on the solid trajectory by applying the Piola transform to the stress
divergence and multiplying the remaining spatial source terms by \(J\).  Using
\eqref{eq:solid_reference_effective_piola_stress}, this transformation gives
%
\begin{align}
	J\sum_\xi\rho_\xi\mbf a_\xi
	= & \;
	\nabla_{\mbf X}\cdot
		\left(
		\mbf{P}^{\prime\prime}
		-B\bar p_E J \mbf F^{-T}
		+J\left(\mbf E\otimes\mbf d\right)\mbf F^{-T}
		\right)
	+J\rho\mbf g
	\notag \\
	&\;
	-J\phi
	\left.
	\p{\gamma}{\theta_{\mc F}}
	\right\vert_{\{S_f\}_{f\in\mc F},\mbf h}
	\mbf F^{-T}\nabla_{\mbf X}\theta_{\mc F}
	\notag \\
	&\;
	-J\phi\mbf F^{-T}
	\left(\nabla_{\mbf X}\mbf h\right)^T
	\left.
	\p{\gamma}{\mbf h}
	\right\vert_{\{S_f\}_{f\in\mc F},\theta_{\mc F}}
	\notag \\
	&\;
	+J\phi\mbf F^{-T}
	\sum_{f\in\mc F}
	\left(
	\bar p_f-\bar p_E-\gamma_f-\omega_f^+
	\right)
	\nabla_{\mbf X}S_f
	\notag \\
	&\;
	+J\sum_\xi\sum_\alpha\dot{c}_\xi^\alpha
	\left(
		\mbf F^{-T}\nabla_{\mbf X}\tau-\mbf v_\xi
	\right).
	\label{eq:solid_reference_overall_momentum}
\end{align}
%
Here
\(\nabla_{\mbf x}\tau=\mbf F^{-T}\nabla_{\mbf X}\tau\), so every gradient in
the conversion-insertion force is evaluated from the solid reference
configuration.  The same pull-back gives
\(\nabla_{\mbf x}\theta_{\mc F}
=\mbf F^{-T}\nabla_{\mbf X}\theta_{\mc F}\) in the thermocapillary force and
\(\nabla_{\mbf x}\mbf h
=\nabla_{\mbf X}\mbf h\,\mbf F^{-1}\) in the capillary-history force.
The same scalar-gradient pull-back carries the dynamic pressure-lag term through
\(\nabla_{\mbf x}S_f=\mbf F^{-T}\nabla_{\mbf X}S_f\).

If a quasi-static solid-skeleton solve is desired, the inertial term on the left
is dropped.  If the fluid velocities are eliminated with
\eqref{eq:relative_flux_velocity_reconstruction}, then the same equation is
solved with $\mbf v_f=\mbf v_{\mc S}+\mbf w_f/\rho_f$ for fluid phases and
$\mbf v_{\mc S}$ for the solid skeleton.  This leaves the solid displacement, the
Lagrangian partial masses, and the fluid relative fluxes as the natural fields
for a reference-configuration implementation.

This quasi-static mechanics--fluid block is only part of the coupled model.
A complete solve must additionally retain the solid component balances
\eqref{eq:solid_reference_solid_component_balance}, the two subsystem energy
balances \eqref{eq:MC_energy_balance}, Gauss' law \eqref{eq:gauss_law},
mechanism kinetics, capillary-history and solid internal-variable evolution,
and the remaining constitutive restrictions appropriate to the selected
specialization.
